% Generated by scripts/build-paper.py — edit paper/src/survey.src.md
\section{A proposed evaluation protocol}\label{sec:a-proposed-evaluation-protocol}

No experiment is reported in this survey. The protocol below is proposed so that the open questions of §7 can be answered with comparable numbers.

\textbf{Tasks.} Six compact tracks: intent/semantic classification, pairwise judging, tool and model routing, narrative variable coding, sequential agent control, and open-set detection with an explicit unknown option. Visual inputs form a separate track, and text-only systems receive the same extracted features.

\textbf{Systems and baselines.} A majority/rule baseline; TF-IDF or embedding plus a linear classifier; SetFit, ModernBERT and GLiClass \citep{arxiv2209_11055,arxiv2412_13663,arxiv2508_07662}; frozen-LM option logits with and without permutation calibration; constrained JSON generation with the same backbone \citep{arxiv2411_15100}; open decision models; version-pinned hosted Jev; a strong LLM; and two cascades (cheap generative, Jev-gated). Supervised systems report their training data and cost.

\textbf{Metrics.} Accuracy or macro-F1 (invalid outputs counted as errors); NLL, Brier and ECE with a stated binning and bootstrap intervals, keeping Choice and Noul separate; risk–coverage curves, AURC and coverage at fixed risk, with thresholds and temperatures fitted on validation data only; flip rates under order permutation, name–rubric rebinding, neutral and random names, negation, a missing correct option and language variants, each against a test–retest floor; single-request p50/p95, amortized per-question time, throughput and end-to-end time as separate columns; API bills per correct decision, and hardware and energy reported separately.

\textbf{Splits and records.} Group splits and bootstraps by scenario or document, so that 577 turn-level decisions are never treated as 577 independent cases. For each item keep the input, schema, option order, full distribution, model ID and version, timestamp, errors, retries and cost. Register the analysis plan before running it.

\textbf{First experiment.} Name–rubric binding × native calibration × selective escalation on a fixed test set. Measure the base error under each binding. Then test whether low-confidence escalation catches the binding-induced errors, and at what cost. This directly asks whether structured, probabilistic errors are in fact easier to intercept (F1 × F2 × F3).
